Votes cast by citizens residing abroad are tracked at the level of voter
nationality, further broken down by host country where the official source
reports it at that granularity.
Except for Spain and Switzerland (see below), 
abroad entries are coded \texttt{CC\_CCZZxx}, where
\texttt{CC} is the country code of the voter nationality and \texttt{xx} is an
country code of the host country
(e.g.\ \texttt{BG\_BGZZDE} = Bulgarian voters in Germany,
\texttt{FR\_FRZZUS} = French voters in the United States).
Where the source does not break results down by host country, numerical
codes are used in place of \texttt{xx}. The special code \texttt{99}
indicates a country-level aggregate of abroad voters without host-country breakdown,
while smaller codes (e.g.\ \texttt{01}, \texttt{02}) are used for multiple host-country 
aggregates where the source provides some breakdown but not a full list of host countries.

\paragraph{Ships at Sea.} Polish sailors voting at sea are also included in this category, as they cast
ballots outside the country and cannot be attributed to a domestic commune.
Their votes are aggregated into a single row (\texttt{PL\_PLZZ98},
\textit{Statki}).

\paragraph{Region-level aggregates.}
Three countries report abroad votes broken down by the voters' \emph{domestic
region of origin} rather than by country of residence.

\textbf{Belgium --- foreign affairs cantons.}
Belgian citizens abroad are counted against the province where they are enrolled,
following the structure of Belgium's overseas electoral census.
These are referred to as `foreign affairs cantons' (\emph{canton affaires étrangères} /
\emph{speciaal kanton Buza} / \emph{Kanton Auswärtige Angelegenheiten})
and carry a five-digit code of the form \texttt{BE\_KAENNN}.
The \bluevar{parent} field is set to the NUTS~3 code of that province
(e.g.\ \texttt{BE21} for Antwerp), if available, or to the country code \texttt{BE} otherwise.

\textbf{Spain --- C.E.R.E.\ (Censo Electoral de Residentes en el Extranjero).}
Spanish citizens registered abroad are counted against the province where they
are enrolled, following the structure of Spain's overseas electoral census.
Each entry carries a five-digit code (\texttt{ES\_NNNNN}, e.g.\
\texttt{ES\_00028}) whose last two digits correspond to the province number
(\texttt{01}--\texttt{52}).
The \bluevar{parent} field is set to the NUTS~3 code of that province
(e.g.\ \texttt{ES300} for Madrid).

\textbf{Switzerland --- cantonal abroad aggregates.}
Swiss cantons with a significant diaspora report separate counts for voters
residing abroad (\textit{Auslandschweizer}).
These entries carry codes in the \texttt{CH\_9xxx} or \texttt{CH\_19xxx} range
(e.g.\ \texttt{CH\_9010} = Zurich abroad voters).
The \bluevar{parent} field is set to the NUTS~3 code of the corresponding
canton (e.g.\ \texttt{CH040} for Zurich).

\paragraph{Other cases.} The following matrix shows which voter-nationality / host-country combinations
have dedicated rows in the dataset.
Columns are voter nationalities; rows are host countries.
A filled circle ($\bullet$) indicates that at least one election file contains
a row for that combination.
Shaded rows correspond to countries that are themselves primary units of
BLUE\_DB.
